\section{The exact threshold at arbitrary finite formal height}
\label{sec:height}

The interpolation proof has an elliptic analogue. This matters at a
supersingular prime, where torsion in the formal parameter has a
different distribution from roots of unity. Let $K$ be a finite
unramified extension of $\mathbb Q_p$, let $F$ be a one-dimensional
commutative formal group law over $\mathcal O_K$, and assume its
reduction has finite height $d$. Set $q=p^d$. We use a fixed integral
formal coordinate, work on the open unit disc of $\mathbb C_p$, normalize
$v(p)=1$, and denote multiplication by $p^k$ by $[p^k]$. In this
section $F[p^j]$ denotes its geometric torsion in this disc.

\paragraph{Theorem (finite height).}
Put $s_j=1/((q-1)q^{j-1})$ for $j\ge1$. Fix $h_0\ge1$,
$\rho\ge s_{h_0}$ and $B\ge0$. For every nonzero
$g\in\mathcal O_{\mathbb C_p}[T]$ with $G_\rho(g)\le B$ there is a
point $z$ of exact order $p^h$ with $v(g(z))\le G_\rho(g)$ and
\[
          h_0\le h\le h_0+\lfloor q^{h_0-1}B\rfloor.
\]
The upper exponent is the smallest possible universal choice.
One entire coset of $F[p^{h_0-1}]$ consists of such points of the same
order. These assertions hold also for nonzero power series with integral
coefficients. The coefficients of the tested function can be ramified
even though the formal group has an unramified field of definition.

\paragraph{The torsion metric.}
In $[p](T)$ the coefficient of $T$ is $p$, the coefficients of
degrees $2,\ldots,q-1$ belong to $p\mathcal O_K$, and the coefficient of
$T^q$ is a unit. The last two facts follow from the definition of height
and the unramified base. For $t=v(z)>0$, the unique lowest term has
valuation $qt$ if $t<s_1$, and $1+t$ if $t>s_1$. In the latter
region repeated application increases the valuation by one and cannot
give zero. At $t=s_1$, the image has valuation at least $1+s_1$.
Thus a nonzero point killed by $p$ has valuation exactly $s_1$, and
induction shows a point of exact order $p^j$ has valuation $s_j$.
The first unit coefficient of $[p^j](T)$ occurs in degree $q^j$ by
composition after reduction. Weierstrass preparation gives $q^j$ roots
in the open disc counted with multiplicity. They are simple, since
translation by a torsion point identifies the local zero with the
zero at the origin, where the derivative is $p^j\ne0$.
It follows that there are $(q-1)q^{j-1}$ points of exact order $p^j$.
The series $F(X,[-1](Y))$ is $(X-Y)$ times a series with unit constant
term. Hence $v(x-y)=v(x\mathbin{\ominus_F}y)$ for torsion parameters.
An integral change of formal coordinate with unit linear coefficient
preserves these distances. For the elliptic facts about height see
Boxall and Grant, \emph{Singular torsion points on elliptic curves},
\S4, Lemma 4.1.

\paragraph{Interpolation and the matching obstruction.}
Write $k=h_0-1$, $L=H-k$, and $M=q^k$, with $H\ge h_0$.
The map $[p^k]:F[p^H]\longrightarrow U=F[p^L]$ is onto by the kernel
count above. The allowed points form the inverse image of
$U\setminus\{0\}$. Choose one lift $x_u$ for each nonzero $u\in U$.
A lift of a point of order $p^j$ has order $p^{j+k}$. For different
nonzero $u,w$ this gives
\[
 v(x_u)=v(u)/M,\qquad v(x_u-x_w)=v(u\mathbin{\ominus_F}w)/M.
\]
Each shell in $U$ has total valuation one. Translation permutes $U$,
so for $P(T)=\prod_{u\ne0}(T-x_u)$ we obtain
\[
 G_\rho(P)=L/M,\qquad v(P'(x_u))=L/M-v(x_u).
\]
Indeed the nonzero differences from $u$, after omitting the index
$0$, range over all nonzero group elements except $u$. Since every
$v(x_u)\le\rho$, every Lagrange basis polynomial
$P(T)/((T-x_u)P'(x_u))$ has Gauss valuation zero.

Divide $g$ by $P$. The quotient is integral because $P$ is monic and
integral. If $G_\rho(g)<L/M$, the remainder has the same Gauss
valuation as $g$. Lagrange interpolation then gives
$\min_{u\ne0}v(g(x_u))\le G_\rho(g)$. If every fibre contained a point
where the inequality failed, choosing those points as lifts would
contradict this conclusion. For a bounded integral series, truncate at
degree $D$ with $(D+1)s_H>G_\rho(g)$ and $(D+1)\rho>G_\rho(g)$.
The tail changes neither the Gauss valuation nor the comparison at
any of the finitely many tested points.

Conversely use $g=P$. At a different lift $z$ of the same $u$,
all factors except $T-x_u$ retain their valuations at $x_u$.
For $k>0$, the difference $z\mathbin{\ominus_F}x_u$ is a nonzero element
of $F[p^k]$ and has valuation at least $s_k>s_{k+1}$, larger than
all the inter-fibre valuations. Therefore
$v(P(z))=L/M-v(x_u)+v(z-x_u)>L/M$. At $x_u$ the value is zero.
For $k=0$ every tested point is a root already.
Consequently a full interval of $L$ shells can be blocked at cost
exactly $L/M$. The first integer $L$ with $L/M>B$ gives the asserted
floor and sharp endpoint. For several independent formal coordinates of
the same height the coordinatewise Gauss-minimizer argument gives the
same bound on each exponent; the one-variable example proves sharpness.

\paragraph{A supersingular example.}
The smooth reduction of $E:y^2=x^3-x$ at $3$ has four
$\mathbb F_3$-points. Its trace is zero and its formal group has height
two. With $h_0=2$, $B=1$, and $\rho\ge1/72$, one has $q=9$ and
$H_{\rm opt}=11$. A whole coset of nine torsion points works, with
order between $3^2$ and $3^{11}$. At $H=10$ the extremal polynomial
uses one lift for each of the $9^9-1$ nonzero points of $F[3^9]$.
The comparison concerns functions integral in the local parameter at the
identity. Applying it to a global rational function requires checking that
its expansion has the stipulated integrality and Gauss bound.

\paragraph{Exact tradeoff between degree and valuation budget.}
There is a second sharp parameter that the same interpolation identity
settles at once. Impose $\deg g\le D_0$ as well as $G_\rho(g)\le B$.
A window of $L$ shells works for every such polynomial precisely when
\[
 L>q^{h_0-1}B\qquad\text{or}\qquad q^L-1>D_0 .
\]
Indeed, below the node count the interpolated remainder is $g$ itself,
whatever its Gauss valuation is. When both inequalities fail, $P$
meets both constraints and blocks every point. Equivalently the optimal
upper exponent is
\[
 h_0+\min\bigl(\lfloor q^{h_0-1}B\rfloor,
                    \lfloor\log_q(D_0+1)\rfloor\bigr).
\]
Thus the same polynomial attains the minimal degree and minimal budget of
an obstruction simultaneously. For instance at $p=3$, height one,
$h_0=2$, the degree-26 obstruction in the first theorem is of minimum
degree even if its budget constraint is removed. Degree at most 25
guarantees a good point through exponent 4 at every budget.

\subsection{A valuation cutoff over ramified base fields}
Uniform shell radii were used to express the answer in terms of $d$
and $h_0$. There is an exact replacement when the base $K/\mathbb Q_p$
is ramified. Keep finite reduction height and take $G=F[p^H]$.
For a fixed positive $\rho$, let
\[
 N=\{z\in G:v(z)>\rho\},\qquad
 A_*=\frac{H-\sum_{z\in N,\ z\ne0}v(z)}{|N|}.
\]
Assume $G\setminus N$ is nonempty. Among integral polynomials
that satisfy $v(g(z))>G_\rho(g)$ for every $z\in G\setminus N$, the
minimum possible $G_\rho(g)$ is $A_*$, and the minimum possible degree
is $[G:N]-1$. One polynomial attains both. Below either threshold
there is an entire good coset of $N$.

\emph{Proof.} The formal group law and inverse preserve the condition
$v(z)>\rho$, so $N$ is a subgroup. On each nonidentity coset the
valuation is constant. Adding an element of $N$ changes a parameter by
an amount of valuation greater than $\rho$. For the distinguished
Weierstrass polynomial of $[p^H]$, the product of its nonzero roots
has valuation $H$, by the derivative $p^H$ at zero and the
unit factor of preparation. The total valuation of one representative
from each nonidentity coset is therefore $A_*$. The valuation of
differences between two such representatives depends only on their
difference coset. Translation in $G/N$ now gives
$v(P'(x_u))=A_*-v(x_u)$ for the product $P$ of their linear factors.
The interpolation, arbitrary-representative argument and matching
example above apply verbatim.\hfill$\square$

This formula only needs the multiset of torsion valuations, available
from the Newton polygon; coset counts and the sum over $N$ follow by
counting the slopes. In particular, if the comparison radius includes
all nonzero torsion valuations, $N=\{0\}$ and the obstruction budget
is exactly $H$ over every base. It suffices here to take
$\rho\ge\max_{0\ne z\in F[p]}v(z)$, since multiplication by $p$ strictly
increases the valuation of a nonzero point in the disc until it reaches
zero. This also explains why the closed formula for the unrestricted
height-$d$ order window can change over ramified bases: cutting by
valuation need no longer cut precisely between two orders.
